% BHU PYQ -- Professional Exam Template
\documentclass[11pt,a4paper]{article}

\usepackage[a4paper,top=2.5cm,bottom=2.5cm,left=2.8cm,right=2.8cm,headheight=14pt]{geometry}
\usepackage{amsmath,amssymb}
\usepackage[T1]{fontenc}
\usepackage[utf8]{inputenc}
\usepackage{lmodern}
\usepackage{microtype}
\usepackage{fancyhdr}
\usepackage{enumitem}
\usepackage{listings}
\usepackage{hyperref}
\usepackage{lastpage}
\usepackage{parskip}

\hypersetup{colorlinks=true,urlcolor=black,linkcolor=black,
  pdftitle={BVCA-401: Cloud Computing -- BHU PYQ}}

\pagestyle{fancy}
\fancyhf{}
\renewcommand{\headrulewidth}{0.4pt}
\renewcommand{\footrulewidth}{0.4pt}
\fancyhead[L]{\small   \textit{Banaras Hindu University}}
\fancyhead[R]{\small   \textit{BVCA-401: Cloud Computing}}
\fancyfoot[C]{\small Page  \thepage\ of \pageref{LastPage}}

\lstset{
  basicstyle=\small tfamily,
  frame=single,
  breaklines=true,
  columns=flexible,
  keepspaces=true
}


\newlist{parts}{enumerate}{2}
\setlist[parts,1]{label=(\alph*),leftmargin=*,itemsep=4pt,topsep=3pt}
\setlist[parts,2]{label=(\roman*),leftmargin=*,itemsep=2pt,topsep=2pt}

\newcommand{\pts}[1]{\hfill   \textit{[#1]}}


\begin{document}

\begin{center}
  {\large\bfseries BANARAS HINDU UNIVERSITY}\\[2pt]
  {
\normalsize B.Voc. Semester IV Examination, 2022-23}\\[6pt]
  \rule{0.8\linewidth}{0.5pt}\\[6pt]
  {\large\bfseries Paper: BVCA-401 --- Cloud Computing}\\[4pt]
  {
\normalsize Time: Three Hours \hfill Maximum Marks: 70}
\end{center}

\rule{\linewidth}{0.4pt}

\smallskip



\noindent   \textit{Instructions: Attempt all sections. Symbols have their usual meanings.}

\bigskip




\noindent {\large\bfseries SECTION --- A}
\rule{\linewidth}{0.4pt}\smallskip




\noindent   \textbf{Question 1.} Attempt all parts: \hfill   \textbf{10 $ \times$ 1 = 10 Marks}

\begin{parts}
  \item Who are the Cloud service providers in a cloud ecosystem?
  \item Who are the Cloud Consumers in a cloud ecosystem?
  \item Who are the Direct customers in a cloud ecosystem?
  \item What is the cloud?
  \item Name the main cloud service models.
  \item Who is considered the father of cloud computing?
  \item Give a few examples of hybrid cloud.
  \item Name some open-source cloud computing platforms.
  \item Give a few examples of public cloud.
  \item Give a few examples of private cloud.
\end{parts}

\bigskip


\noindent {\large\bfseries SECTION --- B}
\rule{\linewidth}{0.4pt}\smallskip




\noindent   \textbf{Question 2.} Attempt any six parts (Answer in about 200 words each): \hfill   \textbf{6 $ \times$ 6 = 36 Marks}

\begin{parts}
  \item What are the advantages of using cloud computing?
  \item What are the characteristics of cloud architecture that separate it from a traditional one?
  \item What does Infrastructure-as-a-Service (IaaS) refer to?
  \item What does the acronym SaaS mean? How does it relate to cloud computing?
  \item What do you know about the Multi-cloud strategy?
  \item Explain different types of cloud computing services.
  \item What is AWS? What types of services does it provide?
  \item What are the basic requirements of a secure cloud software?
  \item Discuss the use of a hypervisor in cloud computing.
\end{parts}

\bigskip


\noindent {\large\bfseries SECTION --- C}
\rule{\linewidth}{0.4pt}\smallskip

Long Answer questions. Attempt any two questions: \hfill   \textbf{2 $ \times$ 12 = 24 Marks}

\medskip



\noindent   \textbf{Question 3.} What is Virtualization in Cloud Computing? Explain its types and benefits.

\medskip



\noindent   \textbf{Question 4.} What are some of the popularly used cloud computing services?

\medskip



\noindent   \textbf{Question 5.} What are the different types of Clouds in Cloud Computing? Explain with examples.

\medskip



\noindent   \textbf{Question 6.} Explain the role of cloud technologies for social networking. Discuss the operational and economic benefits of SaaS.

\vfill
\rule{\linewidth}{0.4pt}

\begin{center}\small
\href{https://bhu-exam.vercel.app/}{Click here to visit our website}\\[4pt]\href{https://me-aryan.vercel.app/}{Click here to visit our contributor}\\[4pt]\href{https://quashan-me.vercel.app/}{Click here to visit our contributor}
\end{center}
\end{document}
